
\newcolumntype{L}{>{\raggedright\arraybackslash}}
\newcolumntype{R}{>{\raggedleft\arraybackslash}}
\newcolumntype{C}{>{\centering\arraybackslash}}

\title{Conception et Implémentation d'une Infrastructure de Sécurité pour une Plateforme Événementielle en Ligne Optimisée par l'IA}

\author{Khaled Saidi}

\diplomaName{Diplôme National d'Ingénieur en Sciences Appliquées et Technologiques}
\speciality{Sécurité des Systèmes Informatiques et Réseaux}

\proFramerName{}
\proFramerSpeciality{}

\academicFramerName{Mme Nour Barrani}
\academicFramerSpeciality{Encadrante académique}

\companyName{TEK-UP Tunisie}

\collegeYear{2024 - 2025}

\proSignSentence{}

\academicSignSentence{J'autorise l'étudiant à faire le dépôt de son rapport de stage en vue d'une soutenance.}

\frenchAbstract{Le présent rapport de fin d'études porte sur la conception et l'implémentation d'une infrastructure de sécurité robuste destinée à une plateforme événementielle en ligne dotée de fonctionnalités d'intelligence artificielle. Dans un contexte où la tenue d'événements virtuels connaît une croissance exponentielle, les enjeux liés à la protection des données personnelles, à la confidentialité des communications et à la disponibilité des services deviennent primordiaux. Ce projet propose une architecture orientée microservices, déployée sur infrastructure cloud OVHcloud, intégrant des mécanismes de sécurité multicouches tels que le chiffrement end-to-end, l'authentification multi-facteurs (MFA), le contrôle d'accès basé sur les rôles (RBAC), ainsi que la conformité au Règlement Général sur la Protection des Données (RGPD). L'intelligence artificielle, incarnée par le modèle Whisper de OpenAI, assure la transcription et la traduction audio en temps réel, enrichissant ainsi l'accessibilité et l'expérience utilisateur de la plateforme. L'intégration du progiciel Odoo garantit, quant à elle, une synchronisation fluide des données événementielles et des comptes utilisateurs.}

\frenchAbstractKeywords{cybersécurité, plateforme événementielle, intelligence artificielle, RGPD, microservices, MFA, RBAC, chiffrement, Whisper, OVHcloud}

\englishAbstract{This final year report addresses the design and implementation of a robust security infrastructure for an online event platform enhanced with artificial intelligence capabilities. As virtual events continue to proliferate, protecting personal data, ensuring communication confidentiality, and guaranteeing service availability have become critical imperatives. The proposed solution adopts a microservices architecture deployed on OVHcloud, incorporating multilayer security mechanisms including end-to-end encryption, multi-factor authentication (MFA), role-based access control (RBAC), and full compliance with the General Data Protection Regulation (GDPR). Artificial intelligence, powered by OpenAI's Whisper model, enables real-time audio transcription and translation, thereby enhancing platform accessibility and user experience. The Odoo ERP integration ensures seamless synchronisation of event data and user accounts.}

\englishAbstractKeywords{cybersecurity, event platform, artificial intelligence, GDPR, microservices, MFA, RBAC, encryption, Whisper, OVHcloud}

\setboolean{wantToTypeCompanyAddress}{false}

\companyEmail{contact@tek-up.de}
\companyTel{+216 71 000 000}
\companyFax{+216 71 000 001}
\companyAddressAR{تونس}
\companyAddressFR{Tunis, Tunisie}
